\section{Importing data from diverse file formats}
    
    Suppose a physics experiment records a variable which is of interest to us, such as temperature, pressure, voltage, count rate etc. These variables are recorded (or sampled) at a sequence of time instants. Now this raw record of data produced by an instrument is almost certainly unfit for direct analysis because
    \begin{enumerate}[a)]
    	\item It may contain unreliable readings.
    	\item It may be out of chronological order.
    	\item It may mix standard and non-standard units.
    	\item It may contain gaps.
    \end{enumerate}
    Then we say that such a dataset must be \textit{pre-processed} before it can be visualised or fitted to a theoretical curve.
    
    For the purpose of pre-processing, the dataset must be first stored and archived in the form of a file. Three file formats are widely physics laboratory exports:
    \begin{enumerate}
    	\item Comma-separated values (CSV files)
    	\item Whitespace-delimited plain text (ASCII files)
    	\item Instrument-specific data (DAT files)
    \end{enumerate}
    
    \subsection{CSV files}
    	A comma-separated values (CSV) file is a plain text data format for storing tabular data where the fields (values) of a record are separated by a comma and each record is a line (i.e. newline separated). For example, the following are the contents of a CSV file:
    	\begin{Verbatim}[frame=single, numbers=left]
    id, name, email
    1, John, john.doe@example.com
    2, Jane, janey72@test.org
    	\end{Verbatim}
    	
    	In software that generally deals with tabular data such as a database or a spreadsheet, the contents of the above CSV file is parsed in the form of a table as follows:
    	\begin{table}[h!]
    		\centering
    		\begin{tabular}{|c|c|c|}
    			\hline
    			\textbf{id} & \textbf{name} & \textbf{email} \\ \hline
    			{1} & {John} & {john.doe@example.com} \\ \hline
    			{2} & {Jane} & {janey72@test.org} \\ \hline
    		\end{tabular}
    	\end{table}
    	
    	The first line is used as the table header and the remaining lines are used as the rows of the table.
    	%In CSV files, the data are delimited using commas and usually carry a single header row. Their apparent simplicity is precisely what makes an inconsistent missing-value marker easy to overlook unless the import step is explicitly configured to recognise it.
    	
    	\subsubsection{Steps to create a CSV file in Windows}
    		\begin{enumerate}[I.]
    			\item 
    		\end{enumerate}
    
    \subsection{ASCII files}
    
    \subsection{DAT files}